\section{Exact minima of the balanced purity}\label{sec:lp}

For a pure state $\psi$ of $n$ qubits and $k=\lfloor n/2\rfloor$ write
$\pi_k(\psi)=\binom nk^{-1}\sum_{|S|=k}\tr\rho_S^2$ for the mean balanced-cut purity. Since
$\tr\rho_S^2\ge2^{-k}$, $\pi_k\ge2^{-k}$, with equality iff $\psi$ is absolutely maximally entangled
(AME$(n,2)$), which happens only for $n=2,3,5,6$ (Rains~\cite{Rains} and Scott~\cite{Scott} for $n\ge8$;
Huber--G\"uhne--Siewert~\cite{HGS} for $n=7$). This section shows that for $n\le18$, $n\ne7,13$,
the minimum of $\pi_k$ is the value of the Rains linear programme and is attained by a graph state
(Table~\ref{tab:minima}; $n=18$ is Theorem~\ref{thm:eighteen}).

\begin{table}[htbp]
\centering
\begin{tabular}{@{}cccl@{}}
\toprule
$n$ & $2^{-k}$ & $\min\pi_k$ (the Rains LP value) & attained by \\
\midrule
4 & $1/4$ & $1/3$ & cluster ($P_4$) or Higuchi--Sudbery state\\
5 & $1/4$ & $1/4$ & AME(5,2), e.g.\ $C_5(1)$\\
6 & $1/8$ & $1/8$ & AME(6,2), e.g.\ $C_6(2,3)$\\
8 & $1/16$ & $3/35$ & graph state (Table~\ref{tab:graphs})\\
9 & $1/16$ & $1/14$ & $C_9(3,4)$\\
10 & $1/32$ & $1/21$ & $C_{10}(1,2,4)$ (Remark~\ref{rem:ten})\\
11 & $1/32$ & $1/28$ & graph state (Table~\ref{tab:graphs})\\
12 & $1/64$ & $5/224$ & $C_{12}(3,5,6)$\\
14 & $1/128$ & $331/27456$ & graph state (Theorem~\ref{thm:fourteenfifteen})\\
15 & $1/128$ & $1/104$ & graph state (Theorem~\ref{thm:fourteenfifteen})\\
16 & $1/256$ & $1/156$ & induced subgraph of $C_{17}(1,3,4,5)$\\
17 & $1/256$ & $1/208$ & $C_{17}(1,3,4,5)$\\
18 & $1/512$ & $7/2288$ & graph state (Theorem~\ref{thm:eighteen})\\
\bottomrule
\end{tabular}
\caption{Minimum of $\pi_k$ for $n\le18$ qubits, $n\ne7,13$ (for $n=7$ the minimum $121/917$ is
Corollary~\ref{cor:seven-global-from-fivecopy}). $C_n(D)$ is the circulant graph joining $i$ and $i\pm d$ for $d\in D$.}
\label{tab:minima}
\end{table}

Write $A_j=\sum_{\mathrm{wt}(P)=j}\langle\psi|P|\psi\rangle^2$ for the weight distribution over the
$4^n$ Pauli strings $P$ ($A_0=1$, $\sum_jA_j=2^n$). Expanding $\rho_S$ in Pauli strings gives
\begin{equation}\label{eq:as}
  a_s := \binom ns^{-1}\sum_{|S|=s}\tr\rho_S^2 = \frac{1}{2^s\binom ns}\sum_{j\le s}\binom{n-j}{s-j}A_j,
\end{equation}
and for pure states $a_s=a_{n-s}$. Rains' shadow inequalities~\cite{Rains} state that the coefficients
$S_j$ of $S(x,y)=\sum_jA_j\bigl(\tfrac{x+3y}{2}\bigr)^{n-j}\bigl(\tfrac{y-x}{2}\bigr)^{j}$ are
non-negative. The \emph{Rains LP} minimises $a_k$ over $A\in\mathbb R^{n+1}_{\ge0}$ with $A_0=1$,
$a_s=a_{n-s}$ and $S_j\ge0$.

\begin{theorem}[exact certificate; LP-1, GS-1]\label{thm:lptable}
\begin{enumerate}
  \item For $n\in\{4,5,6,8,9,10,11,12\}$, $\pi_k(\psi)$ is at least the value in Table~\ref{tab:minima} for every pure $n$-qubit state.
  \item For $n\in\{4,5,6,8,9,10,11,12\}$ the graph state of Table~\ref{tab:graphs} has $\pi_k$ equal to
    the LP bound. Hence $\min\pi_k$ equals the LP bound for these $n$.
\end{enumerate}
\end{theorem}

\begin{proof}
(1) For each $n$, LP-1 exhibits rational multipliers $e_s\in\mathbb Q$, $y_i\ge 0$ such that
$a_k - L A_0 - \sum_s e_s(a_s-a_{n-s}) - \sum_i y_iS_i$, as a linear form in $A$, has non-negative
coefficients. On any feasible $A$ the right-hand terms vanish or are non-negative, so $a_k\ge L$.
(Part~(2) shows that $L$ is the LP value.)
(2) For the graph state $|G\rangle$ and a cut $S$, $\tr\rho_S^2 = 2^{-\operatorname{rk}_2\Gamma[S,S^c]}$,
with $\Gamma[S,S^c]$ the biadjacency matrix over $\mathbb F_2$ (Hein--Eisert--Briegel~\cite{HEB}).
GS-1 computes all ranks exactly. \qedhere
\end{proof}

\begin{table}[htbp]
\centering
\begin{tabular}{@{}c>{\raggedright\arraybackslash}p{8cm}l@{}}
\toprule
$n$ & graph & ranks over all $k$-sets $S$\\
\midrule
4 & path $P_4$ (linear cluster) & $1^{2}\,2^{4}$ \\
5 & $C_5(1)$ & $2^{10}$ \\
6 & $C_6(2,3)$ & $3^{20}$ \\
8 & edges 20 21 31 32 41 43 50 51 52 54 60 61 63 64 70 71 76 & $2^6\,3^8\,4^{56}$ \\
9 & $C_9(3,4)$ & $3^{18}\,4^{108}$ \\
10 & $C_{10}(1,2,4)$ & $4^{132}\,5^{120}$ \\
11 & edges 20 21 31 40 41 42 51 53 54 60 61 63 64 70 71 76 81 83 84 85 86 87 92 93 96 98 $\{10\}$0--6 $\{10\}$9 & $4^{66}\,5^{396}$\\
12 & $C_{12}(3,5,6)$ & $4^{6}\,5^{378}\,6^{540}$ \\
\bottomrule
\end{tabular}
\caption{Graph states attaining the LP bound ($r^m$: $m$ cuts of rank $r$). For $n=11$,
``$\{10\}$0--6'' means vertex 10 is joined to $0,\dots,6$.}
\label{tab:graphs}
\end{table}

\begin{remark}[ten qubits]\label{rem:ten}
The $C_{10}(1,2,4)$ graph state is $3$-uniform, with weight distribution
$A=(1,0,0,0,10,72,160,240,285,200,56)$ (GS-2): a $[[10,0,4]]$ stabilizer state. By \eqref{eq:as},
$\pi_5$ depends on $A_4$ and $A_5$ as well as the distance, so not every $[[10,0,4]]$ state attains
$1/21$: the Petersen graph state has $\pi_5=3/56$ (GS-3). Scott's Table~I~\cite{Scott} lists only
the Petersen weight distribution for $n=10$ (value $3/56$, with a caption noting that other
distributions exist); Zha et al.~\cite{Zha14} give $13/336$ as a heuristic ``theoretical
limitation''. The distribution of $C_{10}(1,2,4)$ is the extremal Type~I enumerator $W_{10,1}$ of
Danielsen--Parker~\cite{DP}, so the value $1/21$ is implicit in published data. We have found no earlier
statement that $\min\pi_5=1/21$.
\end{remark}

\begin{theorem}[exact finite certificates; LP-14, LP-15, GS-14, GS-15]
\label{thm:fourteenfifteen}
The minimum mean balanced-cut purity over all pure qubit states is
\[
 \min_{\psi\in(\mathbb C^2)^{\otimes14}}\pi_7(\psi)=\frac{331}{27456},
 \qquad
 \min_{\psi\in(\mathbb C^2)^{\otimes15}}\pi_7(\psi)=\frac1{104}.
\]
Both minima are attained by graph states.
\end{theorem}
\begin{proof}
For each $n=14,15$, the checks LP-14 and LP-15
verify, coefficient by
coefficient in all $n+1$ Pauli-weight coordinates, the exact rational
Rains-LP dual identity
\[
 a_{\lfloor n/2\rfloor}(A)-L A_0
  =\sum_{s< n/2}e_s(a_s-a_{n-s})
      +\sum_{i=0}^n y_i S_i(A)+\sum_{j=0}^n w_j A_j,
       \quad y_i,w_j\ge0,
\]
where the finite lists of rational coefficients are supplied literally
in the script. This proves the universal lower bounds.

For the graph witnesses with symmetric binary adjacency matrices
specified by the hexadecimal row strings in that script, exact
enumeration of the graph stabilizer gives, respectively,
\[
 \begin{array}{c|c|l}
 n & (A_0,\ldots,A_7) & \text{balanced-cut binary rank histogram}\\ \hline
 14 &(1,0,0,0,0,0,161,576)&5^{60}\,6^{1684}\,7^{1688}\\
 15 &(1,0,0,0,0,0,105,540)&5^{35}\,6^{1380}\,7^{5020}.
 \end{array}
\]
The 14- and 15-vertex graphs have 47 and 56 edges. An independent
enumeration of every $7|(n-7)$ binary cut matrix computes its rank
$r_S$, yielding $\operatorname{tr}\rho_S^2=2^{-r_S}$. Averaging the rank
histograms gives exactly $331/27456$ and $1/104$, matching the
certified lower bounds.
\end{proof}

\begin{theorem}[exact finite certificates; GL-1, GL-2]
\label{thm:sixteenseventeen}
The minimum mean balanced purity is $1/156$ for $n=16$ and $1/208$
for $n=17$. It is attained by the graph state $C_{17}(1,3,4,5)$ at
$n=17$, and by its induced graph on vertices $0,\dots,15$ at $n=16$.
\end{theorem}
\begin{proof}
GL-1 verifies exact Rains-LP dual identities with nonnegative rational
shadow and Pauli-weight coefficients, proving the stated universal
lower bounds. GL-2 computes every balanced cut rank over $\mathbb F_2$.
For $n=16$ the ranks are $5^{10}6^{500}7^{6680}8^{5680}$;
for $n=17$ they are $6^{510}7^{4080}8^{19720}$.
Inserting $2^{-r}$ for each cut of rank $r$ yields $1/156$ and
$1/208$ respectively, matching the lower bounds.
\end{proof}
